% !TEX program = xelatex
\documentclass[10pt,a4paper]{article}
\usepackage[a4paper,margin=22mm,headheight=15pt]{geometry}
\usepackage{fontspec,microtype}
\usepackage{booktabs,longtable,tabularx,array,enumitem,xcolor,hyperref,fancyhdr,lastpage}
\usepackage{amsmath,amssymb,graphicx,titlesec,seqsplit}
\IfFontExistsTF{Times New Roman}{\setmainfont{Times New Roman}}{\setmainfont{Latin Modern Roman}}
\IfFontExistsTF{TeX Gyre Heros}{\setsansfont{TeX Gyre Heros}}{\setsansfont{Arial}}
\definecolor{navy}{HTML}{173458}\definecolor{blue}{HTML}{3366CC}\definecolor{darkgray}{HTML}{44546A}
\hypersetup{colorlinks=true,linkcolor=navy,urlcolor=blue,pdfauthor={RAI Risk Taxonomy 2.0 Project},pdftitle={RAI Risk Taxonomy 2.0 Rebuild Technical Report, English}}
\pagestyle{fancy}\fancyhf{}\lhead{\small RAI Risk Taxonomy 2.0 Master}\rhead{\small Technical Report, English}\cfoot{\small \thepage\ / \pageref{LastPage}}
\titleformat{\section}{\Large\bfseries\color{navy}}{\thesection}{0.7em}{}\titleformat{\subsection}{\large\bfseries\color{darkgray}}{\thesubsection}{0.7em}{}
\setlist[itemize]{leftmargin=1.5em,itemsep=2pt}\setlist[enumerate]{leftmargin=1.7em,itemsep=3pt}\renewcommand{\arraystretch}{1.22}
\newcommand{\hashval}[1]{\texttt{\small\seqsplit{#1}}}\graphicspath{{../output/figures/}}
\begin{document}

\begin{titlepage}\pagecolor{navy}\color{white}\vspace*{25mm}
{\sffamily\bfseries\fontsize{27}{33}\selectfont RAI Risk Taxonomy 2.0\par}\vspace{5mm}
{\sffamily\fontsize{21}{27}\selectfont Rebuild Technical Report\par}\vspace{12mm}
{\Large Risk-eligibility cleaning, EM-based L3 mapping, and master release\par}\vfill
{\large Reference date: 26 August 2026\par}\vspace{3mm}{\normalsize English edition\par}\vspace{18mm}
\end{titlepage}\nopagecolor\color{black}\tableofcontents\newpage

\section{Executive summary}
This report documents the complete rebuild of RAI Risk Taxonomy 2.0 after the initial semantic pipeline assigned an emergency and disaster response context to Violence. The 892 source L4 records were frozen, then subjected to risk-eligibility review, deletion, consolidation, splitting, bilingual rewriting, and peer review before any semantic mapping. The resulting 808 cards satisfy \(892-65-20+1=808\).

EM assigned 612 cards (75.7\%) to existing L3 categories. The remaining 196 cards (24.3\%) were routed to domain-specific Others queues for human decision because of ambiguity between eligible current L3 candidates, weak anchor support, instability, or bilingual disagreement. The final distribution is 599 General, 90 Agentic, and 119 Physical. All 46 source L3 rows remain unchanged, with three derived Others routing rows. All 50 internal validation checks passed.

\section{Failure analysis and redesign}
The earlier procedure assumed that every input was a risk. Terms such as life, safety, emergency, CPR, first aid, and actions increased similarity to Violence even though the card described a normal application context rather than an AI interaction risk. The failure demonstrates that semantic proximity cannot establish risk eligibility and that lexical proximity does not identify the causal harm mechanism.

The revised sequence is: freeze sources and L3; test AI involvement, mechanism, adverse outcome, and affected party; execute explicit deletion, merge, and split instructions; write bilingual L4-level titles and definitions; confirm L1 before L3; compute a provisional L3 drafting anchor; compare each card with the current L3 names and definitions; discard non-risks and cards outside the taxonomy's scope; retain eligible risks without an exact L3 in the confirmed L1's Others queue; require each retained Korean and English definition to name an AI technology; remove formulaic AI involvement modifiers from titles while retaining terms that identify a technical object; review bilingual high-similarity pairs against target and mechanism distinctiveness; extract three representative concepts in each language after deduplication; refit the final constrained EM; and preserve uncertain records in Others with an HD reason and two review candidates. Titles follow the nominal style of the L3 master. Risk or hazard suffixes are retained only where they are semantically necessary.

\section{Data integrity baseline}
\begin{table}[htbp]\centering\caption{Source and final L4 counts}
\begin{tabular}{lrrr}\toprule Domain & Source & Final & Net change\\\midrule
General AI & 591 & 599 & +8\\ Agentic AI & 83 & 90 & +7\\ Physical AI & 218 & 119 & -99\\\midrule Total & 892 & 808 & -84\\\bottomrule\end{tabular}\end{table}

\begin{longtable}{p{0.28\textwidth}p{0.64\textwidth}}\caption{Core source hashes}\\\toprule Source & SHA-256\\\midrule\endfirsthead\toprule Source & SHA-256\\\midrule\endhead
L3 master CSV & \hashval{b154be6ced881cc10d93779447afe3b636bcaeaa6094c399fc4d49076d3812c4}\\ General L4 CSV & \hashval{97a57589bb218a89372bf703cb569c7e5344ff09155dfb74bc624b9a0bdeebe8}\\ Agentic L4 CSV & \hashval{d790d6feda006e0026d582aa0ad0b3d1e8442257eb563583e6077039ab4623cf}\\ Physical L4 CSV & \hashval{d578772484243a6888a2e41afbefb580385debbac2ca70c29c67a1fcebdf349c}\\\bottomrule\end{longtable}

The source-defined fields of all 46 L3 rows were compared cell by cell. The only additional hierarchy rows are G\_Others, A\_Others, and P\_Others, which are derived HD routes rather than formal L3 modifications.

\section{L4 cleaning and peer review}
The archive contains 39 explicit deletions, eight application-context records rejected by the risk-eligibility gate, and one peer-reviewed drop. Fourteen consolidation groups absorbed 20 records, and one multi-meaning record was split into two, adding one net record. Of 100 source records labelled `new', eight context-only records were removed and 87 were redefined and reissued.

\begin{figure}[htbp]\centering\includegraphics[width=0.96\textwidth]{cleaning_reconciliation.png}\caption{L4 reconciliation including the risk-eligibility gate}\end{figure}

Claude-assisted review proposed 713 Korean spacing restorations, eight substantive redefinitions, and one drop. The changes were adopted only after comparison with source meaning. RAI4-0568 was further refined to identify developer conduct, data-subject rights, and organisational liability. The revised RAI4-1157 describes cyberattack enablement, so it was routed from Physical to General and given a strong G\_SYS\_SECADV drafting prior without suppressing the second candidate. Mechanical addition of a risk suffix to every Korean title was discontinued. Harm, infringement, hazard, or risk wording remains only where semantically needed.

All 825 records remaining before the L3 scope gate were compared with the immutable L3 names and definitions. The existing definitions of 548 records already named an AI technology and used a complete risk-statement structure. The other 260 records were amended, without replacing their source meaning, to name an AI system, algorithm, agent, robot, humanoid, learning technology, or model in both languages. The gate archived 12 non-risks or records outside the taxonomy's overall scope. It retained 14 eligible risks that belonged to a confirmed L1 but lacked an exact current L3, together with other boundary cases, in that L1's Others queue.

The terminology review normalised 66 titles by removing formulaic AI involvement modifiers or replacing them with a standard risk noun phrase. All 808 retained titles have an audit trail to the immutable L3 master and terminology families used by ISO, NIST, OECD, UNESCO, UNICEF, WHO, UNODC, WIPO, IMF, or the Korean AI Basic Act. Institutional wording was not copied mechanically. Each title was checked for consistency with the harm, failure, or infringement concept of its L3.

The 813 revised cards produced 66 bilingual high-similarity candidate pairs within the same L3. Similarity generated candidates but never deleted a card automatically. The review retained 61 pairs with distinct protected attributes, weapon types, affected targets, or harm mechanisms. It discarded only 5 less representative cards that added no distinct scope, linking each discarded source row to a retained representative.

\begin{table}[htbp]\centering\caption{Semantic near-duplicate review results}
\begin{tabular}{lrr}\toprule Decision & Candidate pairs & Cards discarded\\\midrule
Retained as distinct scope & 61 & 0\\ Less representative duplicate discarded & 5 & 5\\\midrule Total & 66 & 5\\\bottomrule\end{tabular}\end{table}

\begin{figure}[htbp]\centering\includegraphics[width=0.72\textwidth]{semantic_near_duplicate_review.png}\caption{Decisions for bilingual high-similarity candidate pairs}\end{figure}

\begin{figure}[htbp]\centering\includegraphics[width=0.90\textwidth]{definition_grounding_by_domain.png}\caption{Definitions retained or amended after immutable-L3 and AI-technology review}\end{figure}

\section{Routing and identifiers}
The pipeline applied 103 explicit L1 movement instructions, one separately resolved instruction, and one peer-review semantic route. A separate L1-first cross-domain review then checked 30 records before domain-constrained L3 mapping. For 14 records with no exact L3 in the confirmed L1, the pipeline retained them in that L1's Others queue rather than borrowing a superficially similar L3 from another L1. \texttt{facet} and \texttt{act-type} remain L4 attributes but do not enter the embedding. All L4 identifiers were regenerated as \texttt{L3\_ID + sequence}.

\begin{figure}[htbp]\centering\includegraphics[width=0.90\textwidth]{domain_counts_before_after.png}\caption{L4 counts before and after rebuilding and routing}\end{figure}

\section{EM mapping method}
Korean and English texts were embedded separately using one coherent local BAAI/BGE-M3 snapshot. CLS vectors were L2-normalised. For card \(i\) and candidate L3 \(k\),
\begin{equation}s_{ik}=0.5\cos(\mathbf{x}^{ko}_i,\boldsymbol{\mu}^{ko}_k)+0.5\cos(\mathbf{x}^{en}_i,\boldsymbol{\mu}^{en}_k).\end{equation}
The E-step selects the highest eligible L3 within the confirmed L1. The M-step combines the fixed L3 master anchor with the mean of assigned L4 vectors, using anchor weight 4.0. Five initialisations used seeds 20260826 to 20260830.

Each L4 contributes 3 representative concepts in Korean and English. Auditable lexical profiles and exclusion terms derived from the 46 L3 definitions provide a complementary E-step prior. They do not replace the L3 definitions. Sensitive categories require mechanism-specific terms. Curated drafting hints act as strong priors rather than one-candidate constraints, preserving two reviewable candidates. Non-risks and cards semantically outside the taxonomy are excluded before EM. An eligible risk confirmed at L1 but lacking an exact current L3 is retained in that L1's Others queue. Other HD/Others records preserve ambiguity between eligible candidates. The rejection signals include a joint score below 0.560, a hybrid top-two margin below 0.012, stability below 0.80, raw-anchor score below 0.48, and bilingual Top-1 disagreement at a low margin.

\section{Results}
\begin{table}[htbp]\centering\caption{EM and HD results by domain}
\begin{tabular}{lrrrr}\toprule Domain & Final L4 & EM & HD/Others & HD share\\\midrule
General AI & 599 & 451 & 148 & 24.7\%\\ Agentic AI & 90 & 71 & 19 & 21.1\%\\ Physical AI & 119 & 90 & 29 & 24.4\%\\\midrule Total & 808 & 612 & 196 & 24.3\%\\\bottomrule\end{tabular}\end{table}

\begin{figure}[htbp]\centering\includegraphics[width=0.90\textwidth]{mapping_method_by_domain.png}\caption{EM assignments and HD/Others queues by domain}\end{figure}

\begin{table}[htbp]\centering\caption{EM quality diagnostics}
\begin{tabular}{lrrrr}\toprule Domain & Mean top score & Median margin & Mean stability & Bilingual Top-1 agreement\\\midrule
General AI & 0.7605 & 0.0304 & 0.9850 & 62.4\%\\ Agentic AI & 0.7564 & 0.0360 & 0.9911 & 67.8\%\\ Physical AI & 0.7297 & 0.0241 & 0.9966 & 62.2\%\\\bottomrule\end{tabular}\end{table}

\begin{figure}[htbp]\centering\includegraphics[width=0.98\textwidth]{em_quality_diagnostics.png}\caption{Top similarity, margin, and run-stability distributions}\end{figure}
\begin{figure}[htbp]\centering\includegraphics[width=0.99\textwidth]{largest_l3_categories.png}\caption{Largest EM-assigned L3 categories}\end{figure}

The pre-keyword baseline contained 559 EM assignments and 266 HD records. Keyword augmentation and L3-referenced definition review increased the final EM count to 612, a net recovery of 53 records. Wage polarisation was recovered to G\_SOC\_ECON and cybercrime misuse at scale to G\_SYS\_SECADV. The remaining 196 HD records represent explicit uncertainty preservation and must not be interpreted as automatically confirmed L3 assignments.

\begin{figure}[htbp]\centering\includegraphics[width=0.78\textwidth]{em_baseline_comparison.png}\caption{EM and HD counts before and after keyword augmentation}\end{figure}

\section{Three internal validation methods}
\subsection{Source reconciliation and lineage}
The pipeline checked pre/post source hashes, exact preservation of the 46 L3 rows, the cleaning identity, and representation of all 892 source records in the crosswalk. This detects loss, duplication, or L3 alteration.
\subsection{Convergence and initialisation stability}
Five initialisations per domain retained objective values, iteration counts, and card-level agreement. Low agreement contributed to HD routing.
\subsection{Bilingual and semantic-prerequisite checking}
Korean and English Top-1 categories were stored separately. Sensitive-category prerequisites and raw-anchor support were tested alongside bilingual agreement. Normal activities and application contexts were removed before EM by the risk-eligibility gate.

All 50 automated checks passed. They cover source hashes, L3 immutability, count reconciliation, identifier uniqueness, bilingual completeness, explicit AI-technology naming, L3-style definition structure, title terminology evidence, removal of formulaic AI modifiers, near-duplicate decisions and representative-card lineage, the L3 scope gate, valid L3 references, HD-only Others, crosswalk coverage, three representative concepts, two reviewable candidates, and the five-file release structure.

\section{Web-based human review log}
Every L4 card exposes the top two non-Others L3 candidates with their base EM and hybrid scores. Selecting a candidate opens a prefilled GitHub Issue containing the release identifier and review-snapshot hash. A daily workflow accepts only votes matching the current snapshot, retains the latest vote for each reviewer and card, and produces non-binding recommendations when at least three unique reviewers and a strict majority above 50\% are present. It never reassigns a card. Reassignment occurs only after the user explicitly instructs the system to analyse and apply the logs.

\section{Synchronized master release}
The canonical release contains one L1 CSV, one L1/L2/L3 CSV, and three L4 CSVs. The same records generate the web explorer's \texttt{cards.json} and \texttt{hierarchy.json}, public manifests, Validation Record, bilingual reports, figures, download links, and current-release pointer.

\section{Limitations}
The method is anchor-constrained spherical hard EM rather than a full probabilistic mixture-model EM. Thresholds are internal operating rules and have not been calibrated on an external expert-labelled sample. The 196 Others records form a human-decision queue, not a deletion set. Independent validation should use samples stratified by domain, L3, score, margin, and decision route.

\section{Conclusion}
The rebuild converted a lexical-similarity error into an explicit risk-eligibility design requirement. Combining eligibility screening, immutable L3 definitions, bilingual rewriting, conservative EM, and transparent HD isolation produced an auditable 808-card RAI Risk Taxonomy 2.0 master release.
\end{document}
